\begin{table}[htbp]
\centering\scriptsize
\caption{Leading supported predicted-pathway differences between compartments:
the five largest paired mean CLR differences ($|\Delta$CLR$|$) in each direction
among the $270$ supported tests ($q<0.05$) of the $600$-test family. MetaCyc
identifiers are given in the archived results.}
\label{tab:pathways_compartment}
\begin{tabular}{P{5.0cm}rP{5.0cm}r}
\toprule
Pathway & $|\Delta$CLR$|$ & Pathway & $|\Delta$CLR$|$ \\
\midrule
\multicolumn{2}{l}{\textit{Higher in subsurface than surface}} & \multicolumn{2}{l}{\textit{Higher in surface than subsurface}} \\
L-isoleucine biosynthesis IV & 0.22 & L-alanine biosynthesis (superpathway) & 0.36 \\
Incomplete reductive TCA cycle & 0.15 & tRNA processing & 0.35 \\
Adenine and adenosine salvage III & 0.14 & 8-amino-7-oxononanoate biosynthesis I & 0.33 \\
Geranylgeranyl diphosphate biosynthesis II & 0.12 & Biotin biosynthesis I & 0.31 \\
GDP-mannose biosynthesis & 0.09 & TCA cycle VII (acetate producers) & 0.24 \\
\addlinespace
\multicolumn{2}{l}{\textit{Higher in root-adjacent than shallow-subsurface}} & \multicolumn{2}{l}{\textit{Higher in shallow-subsurface than root-adjacent}} \\
Formaldehyde assimilation II (RuMP cycle) & 0.53 & L-isoleucine biosynthesis IV & 0.38 \\
Palmitate biosynthesis II & 0.49 & GDP-mannose biosynthesis & 0.19 \\
L-lysine biosynthesis II & 0.37 & Kdo transfer to lipid IV$_\mathrm{A}$ III & 0.19 \\
TCA cycle VII (acetate producers) & 0.35 & Glycogen degradation I & 0.17 \\
L-alanine biosynthesis (superpathway) & 0.31 & Purine nucleotide \textit{de novo} biosynthesis II & 0.17 \\
\addlinespace
\multicolumn{2}{l}{\textit{Higher in root-adjacent than surface}} & \multicolumn{2}{l}{\textit{Higher in surface than root-adjacent}} \\
Formaldehyde assimilation II (RuMP cycle) & 0.51 & Kdo transfer to lipid IV$_\mathrm{A}$ III & 0.17 \\
1,4-dihydroxy-2-naphthoate biosynthesis I & 0.31 & L-isoleucine biosynthesis IV & 0.16 \\
Phylloquinol biosynthesis (superpathway) & 0.30 & Starch degradation V & 0.16 \\
L-lysine biosynthesis II & 0.28 & CMP-Kdo biosynthesis I & 0.15 \\
Menaquinol-11 biosynthesis (superpathway) & 0.27 & tRNA processing & 0.14 \\
\bottomrule
\end{tabular}
\end{table}

\begin{table}[htbp]
\centering\scriptsize
\caption{Marker-gene sets in the $975$ annotated MAGs of the companion
catalogue. Genomes: percentage carrying at least one KO in the set. Su, Sh, Ro:
mean relative abundance of carrier genomes in surface, shallow-subsurface and root-adjacent libraries. Leading carriers:
phylum with the most carrier genomes (number of genomes); Act., Actinomycetota;
Pse., Pseudomonadota; Bac., Bacillota. KO memberships are archived with the
analysis.}
\label{tab:traits_genomes}
\begin{tabular}{P{4.2cm}rrrrl}
\toprule
Marker set & Genomes (\%) & Su & Sh & Ro & Leading carriers \\
\midrule
CO dehydrogenase (\textit{coxL}) & 38.5 & 43.8 & 45.5 & 40.4 & Act. (170) \\
{[NiFe]}-hydrogenase (\textit{hyaB}, \textit{hhyL}, \textit{hoxH}) & 22.6 & 17.2 & 31.9 & 15.0 & Act. (126) \\
RuBisCO (\textit{rbcL}/\textit{cbbL}) & 24.6 & 18.2 & 24.1 & 18.1 & Act. (90) \\
Photosystem II D1 (\textit{psbA}) & 0.2 & 0.1 & 0.2 & 0.1 & Pse. (1) \\
Nitrogenase (\textit{nifH}) & 1.0 & 0.5 & 0.3 & 0.8 & Pse. (3) \\
Ammonia/methane monooxygenase (\textit{amoA}/\textit{pmoA}) & 3.3 & 3.3 & 4.7 & 1.9 & Act. (16) \\
Nitrite reductase (\textit{nirK}, \textit{nirS}) & 32.4 & 30.6 & 38.8 & 23.1 & Act. (149) \\
Nitrous-oxide reductase (\textit{nosZ}) & 4.1 & 5.1 & 4.0 & 2.2 & Pse. (14) \\
Urease (\textit{ureC}) & 31.2 & 24.0 & 27.9 & 35.3 & Pse. (132) \\
Trehalose, OtsAB (\textit{otsA}, \textit{otsB}) & 72.7 & 85.6 & 89.0 & 83.8 & Act. (336) \\
Trehalose, TreYZ (\textit{treY}, \textit{treZ}) & 70.1 & 74.3 & 79.1 & 76.5 & Act. (258) \\
Ectoine (\textit{ectA}--\textit{ectC}) & 36.4 & 25.9 & 28.0 & 31.3 & Act. (145) \\
Carotenoids (\textit{crtB}) & 49.3 & 60.6 & 52.7 & 53.4 & Pse. (199) \\
Sporulation (\textit{spo0A}) & 7.3 & 5.1 & 4.3 & 7.5 & Bac. (41) \\
Photolyase (\textit{phrB}) & 23.6 & 34.5 & 17.8 & 26.0 & Pse. (71) \\
Excinuclease A (\textit{uvrA}) & 94.6 & 93.2 & 92.0 & 93.8 & Act. (324) \\
Recombinase A (\textit{recA}) & 92.2 & 88.9 & 90.9 & 92.4 & Act. (324) \\
Catalase (\textit{katE}, \textit{katG}) & 77.4 & 86.7 & 76.3 & 84.1 & Act. (266) \\
Superoxide dismutase (\textit{sodA}) & 82.8 & 68.4 & 67.4 & 78.8 & Pse. (238) \\
Dps (\textit{dps}) & 65.1 & 58.2 & 47.3 & 69.5 & Act. (177) \\
\bottomrule
\end{tabular}
\end{table}

\begin{table}[htbp]
\centering\scriptsize
\caption{PICRUSt2 predictions for the marker-gene sets. Fold changes are paired
site-level ratios of the predicted share of KO abundance between compartments
(Su, surface; Sh, shallow-subsurface; Ro, root-adjacent; sign-flip tests,
Benjamini--Hochberg over $60$ tests). Transect $\rho$: Spearman correlation of
site means with the transect coordinate (Benjamini--Hochberg over $20$ tests).
Agreement $\rho$: Spearman correlation between predicted and genome-derived
shares across $119$ matched libraries. $^{*}$\,$q<0.05$. Predicted shares by
compartment are archived with the analysis.}
\label{tab:traits_picrust}
\begin{tabular}{P{4.2cm}rrrrr}
\toprule
Marker set & Sh/Su & Ro/Su & Ro/Sh & Transect $\rho$ & Agreement $\rho$ \\
\midrule
CO dehydrogenase (\textit{coxL}) & 1.09$^{*}$ & 0.78$^{*}$ & 0.74$^{*}$ & $-0.11$ & $0.71$ \\
{[NiFe]}-hydrogenase (\textit{hyaB}, \textit{hhyL}, \textit{hoxH}) & 1.16$^{*}$ & 0.77$^{*}$ & 0.67$^{*}$ & $0.05$ & $0.74$ \\
RuBisCO (\textit{rbcL}/\textit{cbbL}) & 1.01 & 0.81$^{*}$ & 0.82$^{*}$ & $0.26$ & $0.33$ \\
Photosystem II D1 (\textit{psbA}) & 1.25 & 0.90 & 0.81 & $0.05$ & $0.18$ \\
Nitrogenase (\textit{nifH}) & 1.19 & 1.19 & 1.04 & $0.08$ & $0.04$ \\
Ammonia/methane monooxygenase (\textit{amoA}/\textit{pmoA}) & 1.33 & 1.36 & 1.21 & $-0.24$ & $0.41$ \\
Nitrite reductase (\textit{nirK}, \textit{nirS}) & 0.96 & 0.88$^{*}$ & 0.93 & $-0.23$ & $0.40$ \\
Nitrous-oxide reductase (\textit{nosZ}) & 1.01 & 0.95 & 0.95 & $0.20$ & $-0.10$ \\
Urease (\textit{ureC}) & 0.92$^{*}$ & 0.98 & 1.08$^{*}$ & $-0.27$ & $0.26$ \\
Trehalose, OtsAB (\textit{otsA}, \textit{otsB}) & 1.01 & 0.86$^{*}$ & 0.87$^{*}$ & $-0.71^{*}$ & $0.53$ \\
Trehalose, TreYZ (\textit{treY}, \textit{treZ}) & 0.95 & 0.84$^{*}$ & 0.91 & $-0.67^{*}$ & $0.40$ \\
Ectoine (\textit{ectA}--\textit{ectC}) & 1.21$^{*}$ & 1.22$^{*}$ & 1.05 & $0.36^{*}$ & $0.65$ \\
Carotenoids (\textit{crtB}) & 0.97 & 0.92$^{*}$ & 0.94$^{*}$ & $-0.59^{*}$ & $0.13$ \\
Sporulation (\textit{spo0A}) & 0.91 & 1.41$^{*}$ & 1.49$^{*}$ & $0.56^{*}$ & $0.51$ \\
Photolyase (\textit{phrB}) & 1.03 & 0.94$^{*}$ & 0.93$^{*}$ & $-0.27$ & $-0.14$ \\
Excinuclease A (\textit{uvrA}) & 1.03$^{*}$ & 0.99 & 0.96$^{*}$ & $-0.24$ & $-0.21$ \\
Recombinase A (\textit{recA}) & 1.04$^{*}$ & 1.00 & 0.95$^{*}$ & $-0.57^{*}$ & $0.06$ \\
Catalase (\textit{katE}, \textit{katG}) & 0.91$^{*}$ & 1.09$^{*}$ & 1.20$^{*}$ & $-0.16$ & $0.57$ \\
Superoxide dismutase (\textit{sodA}) & 1.00 & 1.06$^{*}$ & 1.06$^{*}$ & $0.70^{*}$ & $0.66$ \\
Dps (\textit{dps}) & 0.97$^{*}$ & 1.07$^{*}$ & 1.10$^{*}$ & $0.43^{*}$ & $0.76$ \\
\bottomrule
\end{tabular}
\end{table}

\begin{table}[htbp]
\centering\scriptsize
\caption{Predicted MetaCyc pathways most strongly correlated with transect
position (Spearman $\rho$ between site-level CLR values and the transect
coordinate; $60$ sites). Of $200$ pathways tested, $92$ were significant
(Benjamini--Hochberg $q<0.05$): $66$ increased and $26$ decreased eastward.
All pathways shown have $q<0.001$.}
\label{tab:pathway_geography}
\begin{tabular}{P{5.2cm}rP{5.2cm}r}
\toprule
\multicolumn{2}{c}{Increasing eastward} & \multicolumn{2}{c}{Decreasing eastward} \\
\cmidrule(lr){1-2}\cmidrule(lr){3-4}
Pathway & $\rho$ & Pathway & $\rho$ \\
\midrule
Queuosine biosynthesis & $0.75$ & \textit{Bifidobacterium} shunt & $-0.72$ \\
Pantothenate and coenzyme A biosynthesis I & $0.73$ & Glycolysis and Entner--Doudoroff (superpathway) & $-0.63$ \\
tRNA charging & $0.67$ & Galactose degradation I (Leloir pathway) & $-0.59$ \\
Phosphopantothenate biosynthesis I & $0.67$ & UDP-glucose-derived O-antigen building blocks (superpathway) & $-0.56$ \\
Pyrimidine deoxyribonucleotide \textit{de novo} biosynthesis I & $0.65$ & Purine nucleotide degradation II (aerobic) & $-0.56$ \\
preQ$_0$ biosynthesis & $0.63$ & TCA cycle VI (obligate autotrophs) & $-0.53$ \\
6-Hydroxymethyl-dihydropterin diphosphate biosynthesis I & $0.63$ & Palmitate biosynthesis II & $-0.53$ \\
Anhydromuropeptide recycling & $0.62$ & L-Glutamate and L-glutamine biosynthesis & $-0.52$ \\
Gondoate biosynthesis (anaerobic) & $0.62$ & Fatty acid $\beta$-oxidation I & $-0.48$ \\
L-Threonine biosynthesis (superpathway) & $0.61$ & Fatty acid salvage & $-0.48$ \\
\bottomrule
\end{tabular}
\end{table}